\section[Groundwater flow and conservative transport]{Groundwater flow and conservative transport}

\subsection[Spatial discretisation and flow problem]{Spatial discretisation and flow problem}

The problem is set up as a simple two-dimensional flow and transport problem with dimensions and properties as defined in Table \ref{franken_flow}. To construct the MODFLOW model that underlies the transport problem proceed as follows:

%\begin{figure}[b]
%\centering
%\includegraphics[width=80mm]{walter.eps}
%\caption[Initial definition of the model grid for PHT3D Exercise 2]{Initial definition of the model grid for PHT3D %Exercise 2}
%\label{fig:mesh_walter}
%\end{figure}


\begin{itemize}

\item In ORTi, create a new model and save it in a new, separate directory.
Remember to use only alphabets (either lower or UPPER case), numbers, 
underscore (\_) and hyphen (\textemdash) in the path.

\item Under {\color{blue}Parameters} $\rightarrow$ {\color{blue}Model} select the {\color{blue}Model} button and specify the attributes of the new model. Select the values {\color{blue}Xsection} for {\color{blue}Dimension}, {\color{blue}unconfined} for {\color{blue}Type} 
and {\color{blue}MODFLOW series} for {\color{blue}Group} 

\item Under {\color{blue}Parameters} $\rightarrow$ {\color{blue}Model} select the {\color{blue}Grid} button and specify the details of the model domain and the spatial discretization in the dialog that opens up. Define a vertical transect that is 2000m long and 20m thick. The model domain
may be discretized into grid cells of 40m length and 1m height. Note, that the model input in ORTi
is grid-independent and that therefore the model discretization can be easily changed 
at any time later. 

\item Select {\color{blue}Parameters} and select the {\color{blue}Time} button.
Set the {\color{blue}Total Simulation Time} to 18250 days and select a time {\color{blue}Step size}
of 365 days. 

\item Go to {\color{blue}Spatial Attributes} $\rightarrow$ {\color{blue}MODFLOW} $\rightarrow$ {\color{blue}LPF} $\rightarrow$ {\color{blue}lpf.8} and set the {\color{blue}Backg.} value of horizontal hydraulic conductivity to 20 $m/day$. Then press {\color{blue}OK} to confirm the value.

\item Go to {\color{blue}Spatial Attributes} $\rightarrow$ {\color{blue}MODFLOW} $\rightarrow$ {\color{blue}LPF} $\rightarrow$ {\color{blue}lpf.9} and set the {\color{blue}Backg.} value of vertical hydraulic conductivity to 1 $m/day$. Then press {\color{blue}OK} to confirm the value.

\item Go to {\color{blue}Spatial Attributes} $\rightarrow$ {\color{blue}MT3DMS} $\rightarrow$ {\color{blue}BTN} $\rightarrow$ {\color{blue}btn.11} and set the {\color{blue}Backg.} value of effective porosity to be 0.25. Then press {\color{blue}OK} to confirm the value.


\item Implement a hydraulic boundary condition that represents the river at the downstream end of this synthetic catchment by defining a specified head boundary at the downstream side of the model. First, define the type of boundary condition by going to {\color{blue}Spatial Attributes} $\rightarrow$ {\color{blue}MODFLOW} $\rightarrow$ {\color{blue}BAS6} $\rightarrow$ {\color{blue}bas.3 Boundary Conditions} and setting the value to 1 (= active cell) for all cells. After that select {\color{blue}Zone} and use the {\color{blue}Zone Tools} to draw a short line near the effluent end. Choose an elevation of z = 8.5 m for the line and draw it from x = 1m to x = 20m. 
The value for the line (=zone) needs to be set to -1. Use, for example,
$river\_bc$ as {\color{blue}Zone Names}. You will see the selected name displayed near 
the location of the line.
  
By setting the boundary condition type to -1 the hydraulic head at the grid cells 
near the defined line will remain unchanged
over the entire simulation period, i.e., it will remain at the value that corresponds to the elevation of
the river water level. In the next step the initial heads need to be defined. 

\item First the define the global initial hydraulic head for the entire domain (Background value) by going to {\color{blue}Spatial Attributes} $\rightarrow$ {\color{blue}MODFLOW} $\rightarrow$ {\color{blue}BAS6} $\rightarrow$ {\color{blue}bas.6 Initial Heads} and set the
 {\color{blue}Backg.} value to 20m. Then press {\color{blue}OK} to confirm the value.  

\item Define the initial hydraulic head at the location of the river by going to {\color{blue}Spatial Attributes} $\rightarrow$ {\color{blue}MODFLOW} $\rightarrow$ {\color{blue}BAS6} $\rightarrow$ {\color{blue}bas.6 Initial Heads}.
Select {\color{blue}Zone} and use the {\color{blue}Zone Tools} to draw again a short line near the effluent end. Choose an elevation of z = 8.5 m for the line and draw it from x = 1m to x = 20m. 
Note, that you can edit the coordinates manually after you have completed the drawing.
The value for the line (=zone) needs to be set to 8.5. Use, for example,
$river\_head$ as {\color{blue}Zone Names}. You will see the selected name and the selected value
displayed near the location of the line.
  

To define the groundwater recharge for this model you will first need to activate the recharge package,
as it is not automatically activated.
This can be done by going to {\color{blue}Add-in} $\rightarrow$ {\color{blue}MODFLOW modules} 
and ticking the box {\color{blue}RCH}. Then tick {\color{blue}OK}.
You will see that you can now go to {\color{blue}Spatial Attributes} $\rightarrow$ {\color{blue}MODFLOW} $\rightarrow$ {\color{blue}RCH} $\rightarrow$ {\color{blue}rch.2 Recharge value} and set the {\color{blue}Backg.} value to 0.001m, corresponding to an annual groundwater recharge of 365mm per year.
Then press {\color{blue}OK} to confirm the value. 

\end{itemize}

\begin{table}[t]
\caption{Flow and transport parameters for the simulation.}
\begin{center}
\begin{tabular}{lc} 
\toprule
Flow simulation                     & steady state \\
Total simulation period\ $ (days)      $         &  18250            \\
Time step length\ $ (days)      $         &  182.5            \\
Model \  length\ $ (m)      $         &  2000            \\
Model \  thickness\ $(m)$                        &  20            \\
Grid \ spacing\ $\Delta x $ $(m)$                &    40            \\
Grid \ spacing\ $\Delta z $ $(m)$                &    1            \\
Porosity\     $            $         &  0.25           \\
Horizontal hydraulic conductivity\ $(m/day)$               &  20         \\ 
Vertical   hydraulic conductivity\ $(m/day)$               &  1         \\ 
Piezometric head downstream boundary (River)\ $(m)        $         &  8.5           \\ 
Longitudinal dispersivity\ $(m)         $         &  2.5          \\ 
Transverse  vertical dispersivity\ $(m)         $         &  0.025          \\ 
\bottomrule
\end{tabular}
\label{franken_flow}
\end{center}
\end{table}

\subsection[Running MODFLOW]{Running MODFLOW}

You have now completed to set up the flow model and you are ready to run MODFLOW.

\begin{itemize}

\item Under {\color{blue}Parameters} $\rightarrow$ {\color{blue}2.Flow} use the {\color{blue}Write Button} to write all the input files required to run MODFLOW

\item Under {\color{blue}Parameters} $\rightarrow$ {\color{blue}2.Flow} use the 
{\color{blue}Red Arrow} button to start the MODFLOW simulation

\item Check the simulated head contours and assess whether the results are plausible.

\item Try, for example to set the head contours manually at 0.25m intervals, starting from 8.5m.

% Olivier, maybe we can practise the use of some more iPHT3D output features 

\end{itemize}

% Here we had a non-reactive model step in the pmwin version of the tutorial

\subsection{Conservative transport}

For the simulation of the solute transport of the mobile components, 
the model parameters that control advective and dispersive transport need to be entered. 
You will proceed as follows:

\begin{itemize}
\item Go to {\color{blue}Parameters} $\rightarrow$ {\color{blue}3.Transport} and select the  {\color{blue}Parameters} button.
\item In the dialog box that subsequently appears, select {\color{blue}ADV} $\rightarrow$ {\color{blue}adv.1- Major data} $\rightarrow$ {\color{blue}MIXELM: Method} $\rightarrow$  {\color{blue}MMOC} as the Solution Scheme. Leave the default parameters and click {\color{blue}Apply} and Close.

\item Select {\color{blue}Spatial Attributes} $\rightarrow$ {\color{blue}MT3DMS} $\rightarrow$ {\color{blue}DSP}, enter {\color{blue}Backg. longitudinal dispersivity} value of 5 m and click {\color{blue}OK}.

\item Select {\color{blue}Spatial Attributes} $\rightarrow$ {\color{blue}MT3DMS} $\rightarrow$ {\color{blue}DSP}, and enter {\color{blue}Backg. value} for {\color{blue}dsp.3 ratio between horizontal transverse dispersivity} and {\color{blue}longitudinal dispersivity} as 0.01. Click {\color{blue}OK}. This will define the actual value of {\color{blue} horizontal transverse dispersivity} to be 0.05 m (= 5 $\times$ 0.01 m).

\item Add a tracer to the recharge water. This can be done by selecting {\color{blue}Spatial Attributes} $\rightarrow$ {\color{blue}MT3DMS} $\rightarrow$ {\color{blue}btn.23 Recharge}. Allocate an initial concentration of 10 mg/L to the recharge water. Click {\color{blue}OK}.

\item Run the transport model by selecting {\color{blue}Parameters} $\rightarrow$ {\color{blue}3.Transport} $\rightarrow$ {\color{blue}Write}. This will write all required MT3DMS input files. Then press the {\color{blue}Red Arrow} button to begin the transport simulation.

\item Visualise the concentration contours under {\color{blue}Results} $\rightarrow$ {\color{blue}3.Transport} $\rightarrow$ {\color{blue}Tracer}. You can change the {\color{blue}Time Step} (under Results) $\rightarrow$ {\color{blue}1.Model} $\rightarrow$ {\color{blue}Tstep} to visualise the results for other simulation times. Also, after selecting a time it is possible to use the mouse wheel to changes times. 

\item Check the simulated tracer concentration contours and assess whether the results are plausible.

\item Remember to save your ORTi file.

\item End of Day 1.

\end{itemize}


